% Created 2024-01-12 vie 10:04
% Intended LaTeX compiler: pdflatex
\documentclass[11pt]{article}
\usepackage[utf8]{inputenc}
\usepackage[T1]{fontenc}
\usepackage{graphicx}
\usepackage{longtable}
\usepackage{wrapfig}
\usepackage{rotating}
\usepackage[normalem]{ulem}
\usepackage{amsmath}
\usepackage{amssymb}
\usepackage{capt-of}
\usepackage{hyperref}
\usepackage{minted}

\usepackage{parskip}
\usepackage[utf8]{inputenc} %% For unicode chars
\usepackage{placeins}
\usepackage[margin=2.50cm]{geometry}
\usepackage[T1]{fontenc}
\usepackage{mathpazo}
\linespread{1.05}
\usepackage[scaled]{helvet}
\usepackage{courier}
\hypersetup{colorlinks=true,linkcolor=blue}
\RequirePackage{fancyvrb}
\DefineVerbatimEnvironment{verbatim}{Verbatim}{fontsize=\small,formatcom = {\color[rgb]{0.5,0,0}}}
\usepackage{mdframed}
\BeforeBeginEnvironment{minted}{\begin{mdframed}}
\AfterEndEnvironment{minted}{\end{mdframed}}
\setcounter{secnumdepth}{3}
\author{José Manuel Sánchez Álvarez}
\date{}
\title{PHP Form Validation}
\hypersetup{
 pdfauthor={José Manuel Sánchez Álvarez},
 pdftitle={PHP Form Validation},
 pdfkeywords={},
 pdfsubject={},
 pdfcreator={Emacs 28.1 (Org mode 9.5.5)}, 
 pdflang={Spanish}}
\begin{document}

\maketitle
\setcounter{tocdepth}{3}
\tableofcontents

\setlength\parindent{10pt}

\section{Form Validation}
\label{sec:org00e6146}
Form validation is a crucial aspect of web development to ensure that
the data collected from users is correct and useful. In PHP, form
validation involves checking the data submitted by the user against
certain criteria before processing it. Here's an overview of how to
perform basic form validation in PHP:

\subsubsection{Basic Steps for PHP Form Validation:}
\label{sec:org43971c3}
\begin{enumerate}
\item \textbf{Check if the Form is Submitted}: First, check whether the form has
been submitted. This can typically be done by checking if the form's
submit button has been clicked.

\begin{minted}[]{php}
<?
if ($_SERVER["REQUEST_METHOD"] == "POST") {
    // Form was submitted
}
\end{minted}

\item \textbf{Sanitize Input Data}: Before validating, it's important to sanitize
the input to prevent security issues like SQL injections and
cross-site scripting (XSS). PHP provides several functions for this,
such as \texttt{trim()}, \texttt{stripslashes()}, and \texttt{htmlspecialchars()}.

\begin{minted}[]{php}
<?
$name = htmlspecialchars(stripslashes(trim($_POST["name"])));
\end{minted}

\item \textbf{Validate Each Field}: Check each input against specific criteria.
For example, if a field must be filled out, check if it's empty; if
an email is expected, check if it's a valid email format.

\begin{minted}[]{php}
<?
if (empty($_POST["email"])) {
    $emailErr = "Email is required";
} else {
    $email = test_input($_POST["email"]);
    if (!filter_var($email, FILTER_VALIDATE_EMAIL)) {
        $emailErr = "Invalid email format";
    }
}
\end{minted}

\item \textbf{Return Error Messages}: If validation fails, return a meaningful
error message to the user. This helps them understand what
corrections are needed.

\item \textbf{Process the Data}: Once all fields are validated, you can proceed
with processing the data (e.g., storing it in a database, sending an
email, etc.).
\end{enumerate}

\subsubsection{Example: Simple PHP Form Validation}
\label{sec:org10e1a86}
Here's a simple example of a PHP form with validation for a name and an
email field:

\begin{minted}[]{php}
<?php
$nameErr = $emailErr = "";
$name = $email = "";

if ($_SERVER["REQUEST_METHOD"] == "POST") {
    if (empty($_POST["name"])) {
        $nameErr = "Name is required";
    } else {
        $name = test_input($_POST["name"]);
        // check if name only contains letters and whitespace
        if (!preg_match("/^[a-zA-Z-' ]*$/", $name)) {
            $nameErr = "Only letters and white space allowed";
        }
    }

    if (empty($_POST["email"])) {
        $emailErr = "Email is required";
    } else {
        $email = test_input($_POST["email"]);
        // check if e-mail address is well-formed
        if (!filter_var($email, FILTER_VALIDATE_EMAIL)) {
            $emailErr = "Invalid email format";
        }
    }
}

function test_input($data) {
    $data = trim($data);
    $data = stripslashes($data);
    $data = htmlspecialchars($data);
    return $data;
}
?>

<form method="post"
  action="<?php echo htmlspecialchars($_SERVER["PHP_SELF"]);?>">
    Name: <input type="text" name="name" value="<?php echo $name;?>">
    <span class="error">* <?php echo $nameErr;?></span>
    <br><br>
    E-mail: <input type="text" name="email"
                               value="<?php echo $email;?>">
    <span class="error">* <?php echo $emailErr;?></span>
    <br><br>
    <input type="submit" name="submit" value="Submit">
</form>
\end{minted}

In this example, the \texttt{test\_input} function sanitizes the data. The name
and email fields are validated for being required, and additional checks
are performed for proper formatting.


\subsubsection{Conclusion}
\label{sec:org41f11b9}
Form validation in PHP is essential for ensuring data integrity and
security. While the above example covers basic validation, real-world
applications may require more complex checks and security measures, such
as validating against a list of allowed values, checking for uniqueness
in a database, or using CAPTCHA for spam prevention.
\end{document}
